% https://ctan.org/pkg/genealogytree
% https://github.com/zerotoc/pdfinlimg

\documentclass[10pt,paper=a4,ngerman,DIV=11,BCOR=0mm,parskip=half,%
{{ if .Config.DocumentOptions -}}
{{ .Config.DocumentOptions }}%
{{- end }}%
]{scrbook}

\usepackage[ngerman]{babel}
\usepackage{fontspec}
\usepackage{microtype,ellipsis}
\usepackage{libertine}
\usepackage{scrlayer-scrpage}

% \usepackage[
%     %margin=1.25cm
%     left=5mm,
%     right=5mm,
%     top=5mm,
%     bottom=10mm,
%     footskip=5mm
%     ,includeheadfoot
%     ,nohead
%     %,showframe % <- to show the page layout
% ]{geometry}

\usepackage[german=quotes]{csquotes}
\usepackage{longtable,tabu,xltabular}
\usepackage[table,dvipsnames,svgnames]{xcolor}
% https://tex.stackexchange.com/a/323318/71157
\usepackage[pdfencoding=auto,colorlinks=false,hidelinks,breaklinks=true]{hyperref}

\usepackage[most]{tcolorbox}
\usepackage[all]{genealogytree}
%\gtrset{language=german-german}
%\gtrset{image prefix=pictures/}
%\usetikzlibrary{backgrounds}
\usepackage{pgfplots}
\pgfplotsset{compat=1.16}

% \colorlet{mainSection}{gray!30}
% \colorlet{subSection}{gray!15}

% color definitions
\colorlet{accent}{ {{- if $.Config.Options.AccentColor }} {{- $.Config.Options.AccentColor -}} {{ else -}} Black!35! {{- end -}} }
\colorlet{relatedLifeDates}{accent}
\colorlet{fromTo}{accent}
\colorlet{nonContainedPerson}{black}
\colorlet{dotfill}{gray!75}
\colorlet{source-link}{gray!75!black}
%\definecolor{MaterialOrange}{named}{orange}

\makeatletter
\newcommand{\mydots}[1]{%
  {\color{dotfill}%\dotfill}%
    % source: https://tex.stackexchange.com/a/85338/71157
    {\leavevmode \cleaders \hb@xt@ .60em{\hss .\hss }\hfill \kern \z@}%
  }
}
\makeatother


% \newlength{\titleColumnWidth}
% \setlength{\titleColumnWidth}{2.5cm}
\newlength{\ageSpacing}
\setlength{\ageSpacing}{0.3125em}

\renewcommand{\arraystretch}{1.25}


\usepackage{pdfpages}

\usepackage{makeidx}
\makeindex

\usepackage{enotez}
\setenotez{
  reset=true,
  counter-format=Alph,
  list-name=Quellen,
  list-heading={\subsection*{#1}},
}

\DeclareInstance{enotez-list}{sources}{list}{list-type = itemize}

%\DeclareInstance{enotez-list}{custom}{paragraph}{
%  heading = {#1},
%  notes-sep = \baselineskip,
%  format    = \footnotefont,
%  number    = \textsuperscript{#1}
%}


\usepackage[inline]{enumitem}
\setlist[enumerate]{%
  leftmargin=*,%
  label=\fbox{\arabic*},%
  nosep,
  itemjoin={%
  {, }%
  },
  %topsep=0pt,%
  %partopsep=0pt,%
  %noitemsep,%
}
\setlist[itemize]{%
  leftmargin=*,%
  nosep,
  itemjoin={%
  {, }%
  },
  %topsep=0pt,%
  %partopsep=0pt,%
  %noitemsep,%
}
%\setlist[description]{style=nextline}
%\makeatletter
\newlist{datalist}{itemize}{1} % description does not work here
\usepackage{calc}
\setlist[datalist]{
  label=(a),
  wide = 0pt,
  widest = {\bfseries{}Standesamtlich},
  leftmargin =*,
  labelsep = 15pt,
  font = \bfseries\scshape,
  itemsep = 2pt,
  parsep = 1pt,
}
\newlist{simplelist}{itemize}{1} % description does not work here
\setlist[simplelist]{
  label={},
  wide = 0pt,
  widest = {},
  leftmargin =*,
  labelsep = 0pt,
  font = {},
  itemsep = 2pt,
  parsep = 1pt,
}

\newlist{customranking}{itemize}{1} % description does not work here
\setlist[customranking]{
  label={1},
  leftmargin=*,
  align=left,
  labelsep=1em,
  %itemindent=-2.9em, % TODO dynamic
  font=\niceunderline,
  itemsep=0em,
  parsep=0em,
}

\newlist{enumerateranking}{enumerate}{1}
\setlist[enumerateranking]{
  leftmargin=*,%
  label=\fbox{\arabic*},%
  nosep,
  align=right,
  labelsep=0.4em,
  widest= {{- $.Config.Options.Stats.Count -}},
}
\newcommand{\rankingValue}[1]{\niceunderline{#1}:\hspace{0.25em}}
\newcommand{\rankingValueRight}[1]{#1}


%\makeatother


% \RedeclareSectionCommand[
%   beforeskip=-\baselineskip,
%   afterskip=1.0\baselineskip]{subsection}

% destroys clickable hyperref links
%\usepackage{titlesec}
%\titlespacing*{\section}{0pt}{4.0ex plus 1ex minus .2ex}{2.3ex plus .2ex}

%\usetikzlibrary{calc}


% \usepackage{varwidth}
\newtcolorbox{databox}[2][]{enhanced,skin=enhancedlast jigsaw,sharp corners,breakable,
  attach boxed title to top left={xshift=-4mm,yshift=-0.5mm},
  fonttitle=\bfseries\sffamily,varwidth boxed title=0.7\linewidth,
  colbacktitle=accent,
  % TODO title with less padding/sep: titlebottom=0pt,
  colframe=black!20,interior style={top color=white,bottom color=white},
  boxed title style={empty,arc=0pt,outer arc=0pt,boxrule=0pt},
  underlay boxed title={
  \fill[accent] (title.north west) -- (title.north east)
  -- +(\tcboxedtitleheight-1mm,-\tcboxedtitleheight+1mm)
  -- ([xshift=4mm,yshift=0.5mm]frame.north east) -- +(0mm,-1mm)
  -- (title.south west) -- cycle;
  \fill[accent!50!black] ([yshift=-0.5mm]frame.north west)
  -- +(-0.4,0) -- +(0,-0.3) -- cycle;
  \fill[accent!50!black] ([yshift=-0.5mm]frame.north east)
  -- +(0,-0.3) -- +(0.4,0) -- cycle; },
  left=3pt,right=3pt,top=3pt,bottom=3pt,
  title={\strut #2},#1,
  subtitle style={
    boxrule=0pt,
    fontupper={\color{black}},
    colback=accent!30!white,
  },
}

% \renewtcolorbox{databox}[2][]{enhanced,skin=enhancedlast jigsaw,sharp corners,breakable,
%   fonttitle=\bfseries\sffamily,
%   %varwidth boxed title=0.7\linewidth,
%   coltitle=black,
%   colbacktitle=accent,
%   colframe=accent,
%   interior style={top color=white,bottom color=white},
%   boxed title style={empty,arc=0pt,outer arc=0pt,boxrule=0pt},
%   title={#2},#1,
% }


% \renewtcolorbox{databox}[2][]{enhanced,
%   colframe=accent,colback=white,colbacktitle=white,
%   fonttitle=\bfseries,coltitle=black,attach boxed title to top center=
%   {yshift=-0.25mm-\tcboxedtitleheight/2,yshifttext=2mm-\tcboxedtitleheight/2},
%   boxrule=0.25mm,
%   boxed title style={boxrule=0.5mm,
%   frame code={ \path[tcb fill frame] ([xshift=-4mm]frame.west)
%   -- (frame.north west) -- (frame.north east) -- ([xshift=4mm]frame.east)
%   -- (frame.south east) -- (frame.south west) -- cycle; },
%   interior code={ \path[tcb fill interior] ([xshift=-2mm]interior.west)
%   -- (interior.north west) -- (interior.north east)
%   -- ([xshift=2mm]interior.east) -- (interior.south east) -- (interior.south west)
%   -- cycle;} },
%   title={#2},#1,
% }


% simple black attempt
% \renewtcolorbox{databox}[2][]{enhanced,skin=enhancedlast jigsaw,sharp corners,breakable,
%   fonttitle=\bfseries\sffamily,
%   %varwidth boxed title=0.7\linewidth,
%   coltitle=black,
%   colbacktitle=white,
%   interior style={top color=white,bottom color=white},
%   boxed title style={empty,arc=0pt,outer arc=0pt,boxrule=0pt},
%   title={\strut #2},#1,
%   %borderline west={2pt}{0pt}{black},
%   subtitle style={
%     boxrule=0pt,
%     %toprule=0pt,
%     fontupper={\color{black}},
%     colback=white,
%   },
%   boxrule=0pt,
%   titlerule=1pt,
% }

% see page 73 in tcolorbox manual


% \newtcolorbox{databoxinner}[2][]{enhanced,skin=tile,sharp corners,breakable,
%   %fonttitle=,
%   coltitle=black,
%   colbacktitle=accent!30!white,
%   interior style={top color=white,bottom color=white},
%   left=6pt,right=6pt,boxsep=0pt,
%   lefttitle=6pt,righttitle=6pt,toptitle=3pt,bottomtitle=1pt,
%   oversize=0pt,
%   % https://tex.stackexchange.com/a/435489/71157
%   title={\strut #2},#1,
% }




\usepackage{fontawesome}


\makeatletter

\newcommand{\pagerefPerson}[1]{%
{{ if $.Config.Options.Debug.Links -}}
(trying ref to person-#1) %
{{- end }}%
\@ifundefined{r@person-#1}{%
}{%
\scalebox{0.8}{\faChevronCircleRight}\,\pageref{person-#1}%
}%
}

\newcommand{\pagerefPersonBrackets}[1]{%
\@ifundefined{r@person-#1}{%
}{%
(\pagerefPerson{#1})%
}%
}

\newcommand{\refPerson}[2]{%
\@ifundefined{r@person-#1}{%
#2
}{%
\hyperref[person-#1]{#2}%
%{#2}\xspace
}%
}

\makeatother

\newcommand{\birthname}[1]{, geb. \birthnameStyle{#1}}
\newcommand{\birthnameStyle}[1]{\textit{#1}}
\newcommand{\nickname}[1]{\normalfont „#1“}
\newcommand{\lastname}[1]{\textbf{#1}}
\newcommand{\firstname}[1]{\normalfont #1}
% TODO \niceunderline instead of \underline, but not working. \protect or something similar required?
\newcommand{\usedFirstname}[1]{\normalfont\niceunderline{#1}\xspace}
%\newcommand{\usedFirstname}[1]{\normalfont{#1}\xspace}
\newcommand{\additionalFirstname}[1]{\normalfont #1}
\newcommand{\aliasname}[1]{\normalfont -- #1}


{{ $allPersons := allPersons .Generations .ExtraPersons }}
{{ $allPersonsMerged := $allPersons.GetWithMergedTreePositionIDs }}

\newcommand{\unknownPerson}{?}



\usepackage{ifthen}
\newcommand{\age}[1]{%
\ifthenelse{\equal{#1}{\string -1}}
{}
{%
{\color{gray}\footnotesize(mit #1~J.)}%
}%
}

% \newcommand{\duration}[1]{%
% \textcolor{lightgray}{#1}
% }


\makeatletter
\gtr@symbol@factory{Unknown}{%
  \begin{gtr@tikzpicture}[line width=2mm,baseline=-13mm,color=gtrsymbol]%
    \path[use as bounding box] (-0.9,-1.8) rectangle (0.9,0.8);
    \draw (0,0) circle (7mm);
    %\draw[white] (270:7mm) -- (270:18mm);
    %\draw[white] (-0.4,-1.3)--(0.4,-1.3);
  \end{gtr@tikzpicture}%
}{0.011}%
\makeatother



\usepackage[normalem]{ulem} % line breaking problems
\usepackage{soul} % own macro problems
\usepackage{soulutf8} % said to fix umlaut problems
% see https://alexwlchan.net/2017/10/latex-underlines/
\usepackage{contour}
\renewcommand{\ULdepth}{1.8pt}
\contourlength{0.8pt}
\newcommand{\niceunderline}[1]{%
  \uline{\phantom{#1}}%
  \llap{\contour[4]{white}{#1}}%
}


\newcommand{\name}[1]{%
{{ if $.Options.WrapNames -}}
  \textcolor{lightgray}{$\lfloor$}
{{- end -}}
#1
{{- if $.Options.WrapNames -}}
  \textcolor{lightgray}{$\rfloor$}%
{{ end -}}
}

\newcommand{\renderGender}[1]{
  \ifthenelse{\equal{#1}{male}}{\gtrsymMale}{}%
  \ifthenelse{\equal{#1}{female}}{\gtrsymFemale}{}%
  \ifthenelse{\equal{#1}{unknown}}{\gtrsymUnknown}{}%
}
\newcommand{\symbolBox}[1]{
  \makebox[2.7mm]{\hfill
  #1
  \hfill}
}
\newcommand{\renderGenderBox}[1]{
  \symbolBox{\renderGender{#1}}
}
\newcommand{\renderGenderBoxLeft}[1]{
  \makebox[2.5mm]{
  \renderGender{#1}
  }
}


{{ if .Config.CustomHead -}}
    {{ .Config.CustomHead }}
{{ end -}}

%\setcounter{secnumdepth}{0}

\title{
{{- if .Config.Title -}}
  {{ .Config.Title }}
{{- else -}}
  Stammbuch
{{- end -}}
}

{{ with .Config.Author -}}
  {{ if . -}}
    \author{ {{- . -}} }
  {{- end }}
{{- end }}

{{ with .Config.GetNow -}}
  {{ if . -}}
    \date{ {{- .Format "02.01.2006" -}} }
  {{- end }}
{{- end }}


% vorher Installation von Manjaro-Paket ttf-google-fonts-git aus AUR
%\setmainfont{PT Serif}
%\setsansfont{Rubik}

% \setmainfont{Caveat}
% \setsansfont{Permanent Marker}
% \setmonofont{PT Mono}

%%\setmainfont{Quicksand}
%%\setsansfont{Quicksand}
%%\setmonofont{PT Mono}


{{ if not $.Config.Options.Fancharts.Hide }}
  % for fancharts (see basic.tex)
  \gtrset{language=german-german}
  \gtrset{image prefix=pictures/}
  \newcommand{\titlename}[1]{#1}
  \newcommand{\middlename}[1]{%
  {\normalfont #1}%
  }
  \newcommand{\surnbirth}[1]{\normalfont{}, geb. #1}
  \newcommand{\alias}[1]{\surn{#1}}
  \renewcommand{\pref}[1]{\bfseries #1}
  \renewcommand{\surn}[1]{\bfseries #1}


  \makeatletter
  \gtr@db@new@event@store{marriagereligious}

  \gtrset{
      %marriageagereligious/.initial=\gtrsymBuried, % no visible effect at all
      /gtr/event prefix/marriagereligious/.initial=\gtrsymChurch, %  Package pgfkeys Error: I do not know the key '/gtr/event prefix/marriageagereligious', to which you passed 'brel', and I am going to ignore it. Perhaps you misspelled it.
      /gtr/dead/.style={},
    }

  \gtrset{language=german}

  \gtr@symbol@factory{Church}{%
    \begin{gtr@tikzpicture}[line width=3mm,color=gtrsymbol]%
      \path[use as bounding box] (0,0) rectangle (2.7,2.1);
      \draw (0.25,0) -- (2.25,0) -- (2.25,1.5) -- (1.8,1.75) -- (1.35,1.5) -- (1.35,1) -- (0.25,1) -- cycle;
      %\draw (1.35,1) -- (1.35,0);
    \end{gtr@tikzpicture}%
  }{0.011}%

  \gtrset{symlang/Church=kirchl. getraut}

  \makeatother


  \makeatletter
  \gtrDeclareDatabaseFormat{full-ages-one-marriage}{}{%
      \gtrPrintName%
      \begin{gtreventlist}%
          \gtr@list@event{birth}\pgfkeysifdefined{/gtr/database/save/age}{ \textcolor{gray}{(\pgfkeysvalueof{/gtr/database/save/age})}}{}%
          \gtr@list@event{baptism}%
          \gtr@list@event{engagement}%
          \gtrifeventdefined{marriage}{%
              \gtr@list@event{marriage}\pgfkeysifdefined{/gtr/database/save/marriageage}{ \textcolor{gray}{(\pgfkeysvalueof{/gtr/database/save/marriageage})}}{}%
          }{%
              \gtr@list@event{marriagereligious}\pgfkeysifdefined{/gtr/database/save/marriageagereligious}{ \textcolor{gray}{(\pgfkeysvalueof{/gtr/database/save/marriageagereligious})}}{}%
          }%
          \gtr@list@event{divorce}%
          \gtr@list@event{floruit}%
          \gtr@list@event{death}\pgfkeysifdefined{/gtr/database/save/deathage}{ \textcolor{gray}{(\pgfkeysvalueof{/gtr/database/save/deathage})}}{}%
          \gtr@list@event{burial}%
      \end{gtreventlist}%
      \gtr@print@infolist%log.
  }

  \gtrDeclareDatabaseFormat{full-ages}{}{%
      \gtrPrintName%
      \begin{gtreventlist}%
          \gtr@list@event{birth}\pgfkeysifdefined{/gtr/database/save/age}{ \textcolor{gray}{(\pgfkeysvalueof{/gtr/database/save/age})}}{}%
          \gtr@list@event{baptism}%
          \gtr@list@event{engagement}%
          \gtr@list@event{marriage}\pgfkeysifdefined{/gtr/database/save/marriageage}{ \textcolor{gray}{(\pgfkeysvalueof{/gtr/database/save/marriageage})}}{}%
          \gtrifeventdefined{marriagereligious}{%
              \gtr@list@event{marriagereligious}\gtrsymBuried \pgfkeysifdefined{/gtr/database/save/marriageagereligious}{ \textcolor{gray}{(\pgfkeysvalueof{/gtr/database/save/marriageagereligious})}}{}%
          }{}%
          \gtr@list@event{divorce}%
          \gtr@list@event{floruit}%
          \gtr@list@event{death}\pgfkeysifdefined{/gtr/database/save/deathage}{ \textcolor{gray}{(\pgfkeysvalueof{/gtr/database/save/deathage})}}{}%
          \gtr@list@event{burial}%
      \end{gtreventlist}%
      \gtr@print@infolist%
  }

  \gtrDeclareDatabaseFormat{full-ages-no-name}{}{%
      \begin{gtreventlist}%
          \gtr@list@event{birth}\pgfkeysifdefined{/gtr/database/save/age}{ \textcolor{gray}{(\pgfkeysvalueof{/gtr/database/save/age})}}{}%
          \gtr@list@event{baptism}%
          \gtr@list@event{engagement}%
          \gtr@list@event{marriage}\pgfkeysifdefined{/gtr/database/save/marriageage}{ \textcolor{gray}{(\pgfkeysvalueof{/gtr/database/save/marriageage})}}{}%
          \gtr@list@event{marriagereligious}\pgfkeysifdefined{/gtr/database/save/marriageagereligious}{ \textcolor{gray}{(\pgfkeysvalueof{/gtr/database/save/marriageagereligious})}}{}%
          \gtr@list@event{divorce}%
          \gtr@list@event{floruit}%
          \gtr@list@event{death}\pgfkeysifdefined{/gtr/database/save/deathage}{ \textcolor{gray}{(\pgfkeysvalueof{/gtr/database/save/deathage})}}{}%
          \gtr@list@event{burial}%
      \end{gtreventlist}%
      \gtr@print@infolist%
  }

  \gtrDeclareDatabaseFormat{fanchart}{}{%
    \gtrPrintName%
    {\color{gray}
      \begin{gtreventlist}%
          \gtr@list@event{birth}%
          \gtrifeventdefined{birth}{}{%
            \gtr@list@event{baptism}%
          }
          \gtrifeventdefined{marriage}{%
              \gtr@list@event{marriage}%
          }{%
              \gtr@list@event{marriagereligious}%
          }%
          \gtr@list@event{death}%
      \end{gtreventlist}%
    }
  }
  \makeatother
{{ end }}

\begin{document}



{{- if .Config.Attribution }}
\pagestyle{scrheadings}
\clearpairofpagestyles
\cfoot{ \footnotesize \textcolor{gray!50!black}{ {{- .Config.Attribution -}} }}
\ofoot{ \thepage }
{{ end -}}

\makeatletter
\begin{titlepage}
	\centering
  ~

  \vfill

  {\Huge\sffamily\bfseries\@title}
  ~

  \vspace{3em}

  ~
  \begin{tcbraster}[
    raster columns=5,
    raster rows=4,
    raster height=.485\pageheight,
    raster valign=center,
    enhanced,
    size=small,
    sharp corners,
    colframe=white,
    colback=white,
    %watermark overzoom=1.0,
    fit algorithm=hybrid*,
    %raster equal height=all
  ]
  {{ with .Generations -}}
    {{ range $i, $gen := . -}}
      {{ if ge $gen.Index 0 -}}
        {{ range $gen.Persons -}}
          {{ with .GetImageFilename -}}
            \tcbox[
              %colframe=gray!80, colback=gray!10,
              boxsep=0mm, boxrule=0mm,
              colback=white,
            ]{
              \includegraphics[keepaspectratio,width=.8\linewidth]{pictures/{{ . -}} }
            }
          {{ end -}}
        {{ end -}}
      {{ end -}}
    {{ end -}}
  {{ end -}}
  \end{tcbraster}

\end{titlepage}
\makeatother

%\end{document}
%\maketitle


% \chapter*{Einleitung}
% \addcontentsline{toc}{chapter}{Einleitung}
% \markboth{Einleitung}{Einleitung}

{{- with .Config.PreContent }}
{{ . }}
{{- end }}

\tableofcontents



{{ with .Generations }}
  {{ range . }}
    \chapter*{
      {{- if ne .Index 0 }}
        {{- if .IsParent }}
          {{- .Index -}}. Vorfahrengeneration --~
        {{- else -}}
          {{- .Index -}}. Nachfahrengeneration --~
        {{- end }}
      {{- end }}
      {{- .Name -}}
    }
    \addcontentsline{toc}{chapter}{
      {{- if ne .Index 0 }}
        {{- if .IsParent }}
          {{- .Index -}}. Vorfahrengeneration --~
        {{- else -}}
          {{- .Index -}}. Nachfahrengeneration --~
        {{- end }}
      {{- end }}
      {{- .Name -}}
    }
    \markboth{ {{- .Name -}} }{
      {{- if ne .Index 0 }}
        {{- if .IsParent }}
          {{- .Index -}}. Vorfahrengeneration
        {{- else -}}
          {{- .Index -}}. Nachfahrengeneration
        {{- end -}}
      {{ end -}}
    }

    {{ with .Persons }}
      {{ range . }}
        {{ renderFamilybookEntry $.Config . $.Lineages $allPersons }}
      {{- end }}
    {{- end }}
  {{- end }}
{{- end }}


{{ with .ExtraPersons }}
  {{ if .Count }}
    \chapter*{Weitere Personen}
    \addcontentsline{toc}{chapter}{Weitere Personen}
    \markboth{Weitere Personen}{Weitere Personen}
    {{ range .GetPersons }}
      {{ renderFamilybookEntry $.Config . $.Lineages $allPersons }}
    {{- end }}
  {{- end }}
{{- end }}


% Empty pages
{{ if gt .EmptyEntries.Count 0 }}
  \chapter*{Neue Erkenntnisse}
  \addcontentsline{toc}{chapter}{Neue Erkenntnisse}
  \markboth{Neue Erkenntnisse}{Neue Erkenntnisse}
  {{ $showEmptyConfig := $.Config.CopyWithOptionsPersonShowEmpty }}
  {{ range .EmptyEntries.GetPersons -}}
    {{ renderFamilybookEntry $showEmptyConfig . $.Lineages $allPersons }}
  {{ end -}}
{{ end -}}


% Fancharts
{{ if not $.Config.Options.Fancharts.Hide }}
  \chapter*{Übersichten}
  \addcontentsline{toc}{chapter}{Übersichten}
  \markboth{Übersichten}{Übersichten}

  {{ range $.FanchartRoots -}}
    {{ if not (isVollSpitzahn .Root $allPersonsMerged) -}}
      {{ if ne .Generation 0 -}}
        \section*{ {{- .Generation -}}. Elterngeneration}
      {{ end -}}
      \begin{center}
        \begin{autosizetikzpicture}{.4\pageheight}{}
          \gtrfanchart[
            fanchart open right = 210,
            fanchart template = spartan,
            %fanchart boundary color = black!20,
            %fanchart boundary width = 2mm,
            fanchart root style = {black!80, fill=black!10, line width = 1.5mm},
            fanchart inner offset = 3mm,
            fanchart outer offset = 1.5mm,
            fanchart text portrait = inwards,
            fanchart text landscape = clockwise,
            fanchart landscape from level = 0,
            %fanchart segment style for levels = {2,4,6,8,10}{fill=red!10},
            fanchart marker style for levels = {1}{black!10},
            fanchart marker style for levels = {2}{black!20},
            fanchart marker style for levels = {3}{black!30},
            fanchart marker style for levels = {4}{black!40},
            fanchart marker style for levels = {5}{black!50},
            fanchart marker style for levels = {6}{black!60},
            fanchart marker style for levels = {7}{black!70},
            autofill parents male female = {{ $.Config.Options.Fancharts.GenerationsPerChart }},
            %complemented phantom,
            database format=fanchart,
            date format=dd.mm.yyyy,
            ]{
              {{ fanchartParentTree .Root $.Config.Options.Fancharts.GenerationsPerChart .Generation $.Config.Options.MaxParentGenerations }}
            }
        \end{autosizetikzpicture}
      \end{center}
      \newpage
    {{ end -}}
  {{ end -}}    
{{ end -}}


% Spitzahnenliste
{{ if not $.Config.Options.HideSpitzahnen }}
  \chapter*{Spitzahnen}
  \addcontentsline{toc}{chapter}{Spitzahnen}
  \markboth{Spitzahnen}{Spitzahnen}

  Von einem Spitzahn ist mindestens ein Elternteil vollständig unbekannt.
  Bei den mit * markierten Vollspitzahnen gibt es keinerlei Information über Eltern oder Vorfahren.

  {{ with .Generations -}}
    \begin{simplelist}
    {{ range $i, $gen := . -}}
      {{ if ge $gen.Index 0 -}}
        {{ range $gen.Persons -}}
          {{ if (isSpitzahn . $allPersonsMerged) }}
            {{ $name := sub "nameFullNoBirth.tex" . | latexify | trim -}}
            \item {{ sub "personLinkGendered.tex" "person" . "allPersons" $allPersons | trim -}}
            {{ if isVollSpitzahn . $allPersonsMerged -}} * {{ end -}}
            \mydots{} {{ $gen.Index }}. Elterngeneration
          {{ end -}}
        {{ end -}}
      {{ end -}}
    {{ end -}}
  {{ end -}}
  \end{simplelist}
{{ end }}


{{ if not $.Config.Options.Groups.Hide }}
  \chapter*{Gruppen}
  \addcontentsline{toc}{chapter}{Gruppen}
  \markboth{Gruppen}{Gruppen}

  {{ $attrs := $allPersons.GetAttributes }}
  {{ if $attrs }}
    \begin{customranking}
    {{ range $attrs }}
      \item[ {{- .Attribute -}} ]\label{attribute- {{- .Attribute | refName -}} }
      {{ range $i, $p := .Persons.GetPersons -}}
        {{- if $i -}}, {{ end -}}
        {{ sub "personLinkAllRef.tex" "person" $p "allPersons" $allPersons }}
      {{ end -}}
    {{ end -}}
    \end{customranking}
  {{ end -}}
{{ end -}}


{{ if not $.Config.Options.Calendar.Hide }}
  \chapter*{Jahreskalender}
  \addcontentsline{toc}{chapter}{Jahreskalender}
  \markboth{Jahreskalender}{Jahreskalender}

  {{ with allPersons .Generations .ExtraPersons }}
    \newpage
    \section*{Geburtstage}
    \addcontentsline{toc}{section}{Geburtstage}
    \markboth{Jahreskalender}{Geburtstage}

    {{ with $.Stats.BirthDays . -}}
      {{ if len .GetItems }}
        \begin{customranking}
          {{ range .GetItems }}
            \item[{{- .Value -}}]
            {{ range $index, $e := .Persons.GetPersons -}}
              {{- if $index }}, {{ end -}}
              {{ sub "personLinkAllRef.tex" "person" $e "allPersons" $allPersons "hidePages" $.Config.Options.Calendar.HidePages }} %
              ( {{- (sparseDate .GetBirth.Date ).Format "2006" -}} )
            {{- end }}
          {{- end }}
        \end{customranking}
      {{- end }}
    {{- end }}

    \newpage
    \section*{Sterbetage}
    \addcontentsline{toc}{section}{Sterbetage}
    \markboth{Jahreskalender}{Sterbetage}

    {{ with $.Stats.DeathDays . -}}
      {{ if len .GetItems }}
        \begin{customranking}
          {{ range .GetItems }}
            \item[{{- .Value -}}]
            {{ range $index, $e := .Persons.GetPersons -}}
              {{- if $index }}, {{ end -}}
              {{ sub "personLinkAllRef.tex" "person" $e "allPersons" $allPersons "hidePages" $.Config.Options.Calendar.HidePages }} %
              ( {{- (sparseDate .GetDeath.Date ).Format "2006" -}} )
            {{- end }}
          {{- end }}
        \end{customranking}
      {{- end }}
    {{- end }}


    {{ if false }}
    \newpage
    \section*{Hochzeitsjahrestage}
    \addcontentsline{toc}{section}{Hochzeitsjahrestage}
    \markboth{Jahreskalender}{Hochzeitsjahrestage}

    {{ with $.Stats.MarriageDays . -}}
      {{ if len .GetItems }}
        \begin{customranking}
          {{ range .GetItems }}
            {{- $value := .Value -}}
            \item[{{- $value.Format "02.01." -}}]
            {{ range $index, $e := .Persons.GetPersons }}
              {{ if $index }} \gtrsymMarried~{{ end -}}
              {{ sub "personLinkAllRef.tex" "person" $e "allPersons" $allPersons "hidePages" $.Config.Options.Calendar.HidePages }} %
            {{- end }} %
            ( {{- $value.Format "2006" -}} )
          {{- end }}
        \end{customranking}
      {{- end }}
    {{- end }}
    {{- end }}% if false
  {{- end }}

{{ end -}}



{{ if not $.Config.Options.Stats.Hide }}
\chapter*{Statistik}
\addcontentsline{toc}{chapter}{Statistik}
\markboth{Statistik}{Statistik}

    {{ with $allPersons }}
      {{ $total := .Count -}}
      {{ $aliveCount := $.Stats.AliveCount . -}}

      {{ if not $.Config.Options.Stats.HideNumbers }}
        \section*{Anzahl Personen im Stammbuch}
        \addcontentsline{toc}{section}{Anzahl Personen im Stammbuch}

        {{ $male := $.Stats.MaleCount . -}}
        {{ $female := $.Stats.FemaleCount . -}}
        {{ $deadCount := $.Stats.DeadCount . -}}
        {{ $marriedCount := $.Stats.MarriedCount . -}}
        \fbox{\textbf{ {{- $total -}} }%
        }

        \begin{itemize}
          \item davon \textbf{ {{- $female -}} } weiblich ({{ percentageInt $female $total 0 }}\%) und \textbf{ {{- $male -}} } männlich ({{ percentageInt $male $total 0 }}\%)
          {{ if $.Config.Options.AllowLiveData -}}
          \item davon \textbf{ {{- $aliveCount -}} } am Leben ({{ percentageInt $aliveCount $total 0 }}\%) und \textbf{ {{- $deadCount -}} } gestorben ({{ percentageInt $deadCount $total 0 }}\%)
          {{ end -}}
          %\item davon \textbf{ {{- $marriedCount -}} } verheiratet
        \end{itemize}
      {{- end }}

      {{ if not $.Config.Options.Stats.HideGenerations }}
        {{ if gt (len $.Generations) 1 }}
          {{ with $.Generations }}
            \section*{Generationen}
            \addcontentsline{toc}{section}{Generationen}

            \begin{itemize}
              {{ range . }}
                {{- if ne .Index 0 }}
                \item \textbf{
                    {{- if .IsParent }}
                      {{- .Index -}}. Elterngeneration --
                    {{- else -}}
                      {{- .Index -}}. Kindergeneration --
                    {{- end -}} }
                  {{ .Name }} \mydots{} {{- len .Persons }} bekannt
                {{- end }}
              {{- end }}
            \end{itemize}
          {{- end }}
        {{- end }}
      {{- end }}

      {{ if not $.Config.Options.Stats.HideAlive }}
        {{ if and $.Config.Options.AllowLiveData (le $aliveCount $.Config.Options.Stats.Count) }}
          \section*{Am Leben}
          \addcontentsline{toc}{section}{Am Leben}

          {{ with $.Stats.AlivePersons . -}}
          {{ if gt (len .GetPersons) 0 -}}
            \begin{enumerateranking}
                {{ range .GetPersons -}}
                  \item {{ sub "personLinkAllRef.tex" "person" . "allPersons" $allPersons }}
                {{- end }}
              \end{enumerateranking}
            {{- end }}
          {{- end }}
        {{ end -}}
      {{ end -}}

      {{ if not $.Config.Options.Stats.HideEarliestBirth }}
        \section*{Frühstes Geburtsdatum}
        \addcontentsline{toc}{section}{Frühestes Geburtsdatum}

        {{ with $.Stats.OldestBirths . -}}
          \begin{enumerateranking}
            {{ range .GetTop $.Config.Options.Stats.Count }}
              \item
              {{ range .Persons.GetPersons -}}
                {{ sub "personLinkAllRef.tex" "person" . "allPersons" $allPersons }}
              {{- end }}
              \mydots{}
              \rankingValueRight{ {{- .Value -}} }
            {{- end }}
          \end{enumerateranking}
        {{- end }}
      {{- end }}

      {{ if not $.Config.Options.Stats.HideLatestBirth }}
        \section*{Letztes Geburtsdatum}
        \addcontentsline{toc}{section}{Letztes Geburtsdatum}

        {{ with $.Stats.NewestBirths . -}}
          \begin{enumerateranking}
            {{ range .GetTop $.Config.Options.Stats.Count }}
              \item
              {{ range .Persons.GetPersons -}}
                {{ sub "personLinkAllRef.tex" "person" . "allPersons" $allPersons }}
              {{- end }}
              \mydots{}
              \rankingValueRight{ {{- .Value -}} }
            {{- end }}
          \end{enumerateranking}
        {{- end }}
      {{- end }}

      {{ if not $.Config.Options.Stats.HideOldest }}
        \section*{Höchstes Alter}
        \addcontentsline{toc}{section}{Höchstes Alter}

        {{ with $.Stats.Oldest . -}}
          {{ $data := .GetTop $.Config.Options.Stats.Count }}
          \begin{enumerateranking}
            {{ range $data }}
              \item 
              {{ range .Persons.GetPersons -}}
                {{ sub "personLinkAllRef.tex" "person" . "allPersons" $allPersons }}
              {{- end }}
              \mydots{}
              \rankingValueRight{ {{- .Value -}} }
            {{- end }}
          \end{enumerateranking}
        {{- end }}
      {{- end }}

      {{ if not $.Config.Options.Stats.HideDeathAge }}
        \section*{Sterbealter}
        \addcontentsline{toc}{section}{Sterbealter}

        {{ with $.Stats.DeathAges . -}}
          {{ if len .GetItems }}
            \begin{center}
              \begin{tikzpicture}
                \begin{axis}[
                  xlabel=Sterbealter,
                  ylabel={Anzahl},
                  enlarge x limits=0.015,
                  legend style={at={(0.5,-0.15)},anchor=north,legend columns=-1},
                  ybar,
                  bar width=1.5pt,
                  %nodes near coords,
                  ymin=0,
                  xmin=0, xmax=110,
                  %symbolic x coords={Januar,Februar,März,April,Mai,Juni,Juli,August,September,Oktober,November,Dezember},
                  ]
                  \addplot coordinates {
                    {{ range .GetItems -}}
                      ({{ .Value }}, {{- .Persons.Count -}})
                    {{ end -}}
                  };
                \end{axis}
              \end{tikzpicture}
            \end{center}
            \begin{enumerateranking}
              {{ range .GetItems }}
                % \textbf{ {{- .Persons.Count -}} }x:
                \item[] \rankingValue{ {{- .Value }} Jahre}

                {\small
                {{- range $index, $e := .Persons.GetPersons -}}
                  {{- if $index }}, {{ end -}}
                  {{ sub "personLinkAllRef.tex" "person" $e "allPersons" $allPersons }}
                {{- end }}
                }
              {{- end }}
            \end{enumerateranking}
          {{ end -}}
        {{ end -}}
      {{ end -}}


      {{ if not $.Config.Options.Stats.HideMarriageAge }}
        \section*{Heiratsalter}
        \addcontentsline{toc}{section}{Heiratsalter}

        \textbf{Durchschnitt}: {{ printf "%.1f" ($.Stats.AvgMarriageAge .) }} Jahre

        {{ with $.Stats.MarriageAges $allPersonsMerged -}}
          {{ if len .GetItems }}
            \begin{center}
              \begin{tikzpicture}
                \begin{axis}[
                  xlabel=Heiratsalter,
                  ylabel={Anzahl},
                  enlarge x limits=0.015,
                  legend style={at={(0.5,-0.15)},anchor=north,legend columns=-1},
                  ybar,
                  bar width=1.5pt,
                  %nodes near coords,
                  ymin=0,
                  xmin=0, xmax=110,
                  %symbolic x coords={Januar,Februar,März,April,Mai,Juni,Juli,August,September,Oktober,November,Dezember},
                  ]
                  \addplot coordinates {
                    {{ range .GetItems -}}
                      ({{ .Value }}, {{- .Persons.Count -}})
                    {{ end -}}
                  };
                \end{axis}
              \end{tikzpicture}
            \end{center}
            \begin{customranking}
              {{ range .GetItems }}
                \item[ {{- .Value }} Jahre] % \textbf{ {{- .Persons.Count -}} x}:
                {{ range $index, $e := .Persons.GetPersons -}}
                  {{- if $index }}, {{ end -}}
                  {{ sub "personLinkAllRef.tex" "person" $e "allPersons" $allPersons }}
                {{- end }}
              {{- end }}
            \end{customranking}
          {{ end -}}
        {{ end -}}
      {{ end -}}


      {{ if not $.Config.Options.Stats.HideTotalAgeLiving }}
        \section*{Gesamtalter aller Lebenden}
        \addcontentsline{toc}{section}{Gesamtalter aller Lebenden}

        {{ with $.Stats.TotalAgeLiving . -}}
          {{ . }} Jahre
        {{- end }}
      {{- end }}


      {{ if not $.Config.Options.Stats.HideLongestMarriages }}
        \section*{Längste Ehen}
        \addcontentsline{toc}{section}{Längste Ehen}

        {{ with $.Stats.LongestMarriages . -}}
          {{ $data := .GetTop $.Config.Options.Stats.Count }}
          {{ if len $data }}
            \begin{enumerateranking}
              {{ range $data }}
                \item
                {{ range $i, $p := .Persons.GetPersons -}}
                  {{ if $i }} \gtrsymMarried{}\\ {{ end -}}
                  {{ with $p -}}
                    {{ sub "personLinkAllRef.tex" "person" . "allPersons" $allPersons }}
                  {{- end }}
                {{- end }}
                \mydots{}
                \rankingValueRight{ {{- .Value -}} }
              {{- end }}
            \end{enumerateranking}
          {{- end }}
        {{- end }}
      {{- end }}


      {{ if not $.Config.Options.Stats.HideLongestName }}
        \section*{Längster Name}
        \addcontentsline{toc}{section}{Längster Name}

        {{ with $.Stats.LongestName . -}}
          \begin{enumerateranking}
            {{ range .GetTop $.Config.Options.Stats.Count }}
              \item
              {{ range .Persons.GetPersons -}}
                {{ sub "personLinkAllRef.tex" "person" . "allPersons" $allPersons }}
              {{- end }}
              \mydots{}
              \rankingValueRight{ {{- .RankingValue }} Buchstaben}
            {{- end }}
          \end{enumerateranking}
        {{- end }}
      {{- end }}


      {{ if not $.Config.Options.Stats.HideShortestName }}
        \section*{Kürzester Name}
        \addcontentsline{toc}{section}{Kürzester Name}

        {{ with $.Stats.ShortestName . -}}
          \begin{enumerateranking}
            {{ range .GetTop $.Config.Options.Stats.Count }}
            \item
            {{ range .Persons.GetPersons -}}
              {{ sub "personLinkAllRef.tex" "person" . "allPersons" $allPersons }}
            {{- end }}
            \mydots{}
            \rankingValueRight{ {{- .RankingValue }} Buchstaben}
          {{- end }}
          \end{enumerateranking}
        {{- end }}
      {{- end }}


      {{ if not $.Config.Options.Stats.HideMostChildren }}
        \section*{Meiste Kinder}
        \addcontentsline{toc}{section}{Meiste Kinder}

        {{ with $.Stats.ChildrenCount . -}}
          \begin{enumerateranking}
            {{ range .GetTop $.Config.Options.Stats.Count }}
              \item
              {{ range .Persons.GetPersons -}}
                {{ sub "personLinkAllRef.tex" "person" . "allPersons" $allPersons }}
              {{- end }}
              \mydots{}
              \rankingValueRight{ {{- .Value }} Kinder}
            {{- end }}
          \end{enumerateranking}
        {{- end }}
      {{- end }}


      {{ if not $.Config.Options.Stats.HideBirthMonths }}
        \section*{Geburtsmonate}
        \addcontentsline{toc}{section}{Geburtsmonate}

        {{ with $.Stats.BirthMonths . -}}
          \begin{center}
            \begin{tikzpicture}
              \begin{axis}[
                xlabel=Monat,
                ylabel={Anzahl Geburten},
                enlarge x limits=0.15,
                legend style={at={(0.5,-0.15)},anchor=north,legend columns=-1},
                ybar,
                bar width=12pt,
                nodes near coords,
                ymin=0,
                xmin=1,
                xmax=12,
                xtick={1,...,12},
                %symbolic x coords={Januar,Februar,März,April,Mai,Juni,Juli,August,September,Oktober,November,Dezember},
                ]
                \addplot coordinates {
                  {{ range .GetItems -}}
                    ({{ .Value }}, {{- .Persons.Count -}})
                  {{ end -}}
                };
              \end{axis}
            \end{tikzpicture}
          \end{center}
          %\begin{enumerateranking}
          %  {{ range .GetItems }}
          %    \item[{{ .Value }}] \textbf{ {{- .Persons.Count -}} }x
          %    ({{ range $index, $e := .Persons.GetPersons -}}
          %      {{- if $index }}, {{ end -}}
          %      \refPerson{ {{- $e.GetBestID -}} }{%
          %        \name{ {{- $e.GetName.GetUsedFirst | latexify -}} } \pagerefPerson{ {{- $e.GetBestID -}} }%
          %      }
          %    {{- end }})
          %  {{- end }})
          %\end{enumerateranking}
        {{- end }}
      {{- end }}


      {{ if not $.Config.Options.Stats.HideMostCommonLastNames }}
        \section*{Häufigste Nachnamen}
        \addcontentsline{toc}{section}{Häufigste Nachnamen}

        {{ with $.Stats.MostCommonLastNames . -}}
          \begin{enumerateranking}
            {{ range .GetTop $.Config.Options.Stats.Count }}
              \item \rankingValue{ {{- .RankingValue }}x \textbf{ {{- .Value -}} }%
              }

              {\small
              {{- range $index, $e := .Persons.GetPersons -}}
                {{- if $index }}, {{ end -}}
                {{ sub "personLinkAllRef.tex" "person" $e "allPersons" $allPersons }}
              {{- end }}
              }
            {{- end }}
          \end{enumerateranking}
        {{- end }}
      {{- end }}


      {{ if not $.Config.Options.Stats.HideMostCommonFirstNames }}
        \section*{Häufigste Vornamen}
        \addcontentsline{toc}{section}{Häufigste Vornamen}

        {{ with $.Stats.MostCommonFirstNames . -}}
          \begin{enumerateranking}
            {{ range .GetTop $.Config.Options.Stats.Count }}
              \item \rankingValue{ {{- .RankingValue }}x \textbf{ {{- .Value -}} }%
              }

              {\small
              {{- range $index, $e := .Persons.GetPersons -}}
                {{- if $index }}, {{ end -}}
                {{ sub "personLinkAllRef.tex" "person" $e "allPersons" $allPersons }}
              {{- end }}
              }
            {{- end }}
          \end{enumerateranking}
        {{- end }}
      {{- end }}


      {{ if not $.Config.Options.Stats.HideUniqueFirstNames }}
        \section*{Einmalige Vornamen}
        \addcontentsline{toc}{section}{Einmalige Vornamen}

        {{ with $.Stats.UniqueFirstNames . -}}
          \begin{enumerateranking}
            {{ range .RankingItems }}
              \item
              {{ range $index, $e := .Persons.GetPersons -}}
                {{- if $index }}, {{ end -}}
                {{ sub "personLinkAllRef.tex" "person" $e "allPersons" $allPersons }}
              {{- end }}
            {{- end }}
          \end{enumerateranking}
        {{- end }}
      {{- end }}


      {{ if not $.Config.Options.Stats.HideCommonBirthdays }}
        \section*{Gemeinsame Geburtstage}
        \addcontentsline{toc}{section}{Gemeinsame Geburtstage}

        {{ with $.Stats.NonUniqueBirthdays . -}}
          {{ if gt (len .RankingItems) 0 }}
          \begin{customranking}
              {{ range .RankingItems }}
              % \textbf{ {{- .Persons.Count -}} }x
              \item[] {{- .RankingValue -}}

                {{ range $index, $e := .Persons.GetPersons -}}
                  {{- if $index }}, {{ end -}}
                  {{ sub "personLinkAllRef.tex" "person" $e "allPersons" $allPersons }}
                {{- end }}
              {{- end }}
            \end{customranking}
          {{- end }}
        {{- end }}
      {{- end }}


      {{ if not $.Config.Options.Stats.HideDuplicateNames }}
        \section*{Mehrfache Namen}
        \addcontentsline{toc}{section}{Mehrfache Namen}

        {{ with $.Stats.NonUniqueNames . -}}
          {{ if gt (len .RankingItems) 0 }}
            \begin{customranking}
              {{ range .RankingItems }}
                % \textbf{ {{- .Persons.Count -}} }x
                \item[ {{- .Value -}}:]
                
                ~\newline
                {\small
                {{- range $index, $e := .Persons.GetPersons -}}
                  {{- if $index }}, {{ end -}}
                  {{ sub "personLinkAllRef.tex" "person" $e "allPersons" $allPersons }}
                {{- end }}
                }
              {{- end }}
            \end{customranking}
          {{- end }}
        {{- end }}
      {{- end }}
  {{- end }}
{{- end }}


{{ with $.Config.Sources -}}
  {{ if len . -}}
    \chapter*{Quellen}
    \addcontentsline{toc}{chapter}{Quellen}
    \markboth{Quellen}{Quellen}
    {{ range . -}}
      {{ $sourceTitle := .ID -}}
      \section*{ {{- $sourceTitle -}} }
      \addcontentsline{toc}{section}{ {{- $sourceTitle -}} }
      \markboth{Quellen}{ {{- $sourceTitle -}} }
      \label{ {{- labelify "src:full" .ID -}} }
      {{ .Description -}}

      {{ if .PDF }}
        \includepdf[pages={-}, fitpaper=false, scale=.9
        {{- if .Rotate -}} ,angle={{ .Rotate }} {{ end -}} ]{ {{- .PDF | expandPath -}} }
      {{ end -}}
    {{ end -}}
  {{ end -}}
{{ end -}}


\clearpage
\addcontentsline{toc}{chapter}{Index}
\raggedright
\printindex


{{- with .Config.PostContent -}}
{{ . }}
{{- end }}


\end{document}
